\section{从联合界到一致收敛}
\label{sec:13-Machine-Learning/10-Learning-Theory/02}

上一节末尾给出了两条修复路线。第一条是物理拆分——留出一批数据专门做评估，让选出的假设对评估集而言重新变回``固定''的。这条路干净，但评估样本不参与训练，在小数据场景下代价不小。

第二条路更激进：继续用同一批数据做选择，但要求 Hoeffding 式的高概率界对假设类中的每一个假设同时成立。无论 ERM 从中挑出了谁，挑中者的泛化误差都已经受控。这就是一致收敛（uniform convergence）的思路——不是让数据依赖的模型重新变独立，而是提前给所有候选同时建立护栏。

两条路的核心差别在哪？拿具体的数字感受一下。

假设你有一个包含 1000 个候选规则的假设类 \(\mathcal{H}\)，10000 个训练样本，你希望以至少 95\% 的置信度控制泛化误差。对单个固定假设，Hoeffding 给出的界宽大约是 0.014——训练误差和真风险的差距以 95\% 概率不超过 1.4 个百分点。但 ERM 输出的 \(\widehat h\) 不是固定的——它是从 1000 个候选里被特意挑出来的、训练误差最小的那个。它的训练误差已经不再是真风险的无偏估计。

问题可以这样想：你有 1000 张彩票，开奖规则是每张彩票的``训练表现''围绕其``真实水平''随机波动。你看了所有 1000 张彩票的这次开奖结果，挑出表现最好的那张，然后声称这张彩票的真实水平就是它刚才的表现。显然，挑出来的那张大概率正是波动最有利的那张——它的训练表现高于它的真实水平。候选越多，这个效应越强。

\subsection{Union bound 一次买断所有坏事件}

0-1 损失下，对每一个固定的 \(h \in \mathcal{H}\)，Hoeffding 不等式控制单个假设的评估误差超出 \(\varepsilon\) 的概率：

\[
A_h = \bigl\{|\widehat R(h) - R(h)| > \varepsilon\bigr\},
\qquad
P(A_h) \le 2e^{-2n\varepsilon^2}.
\]

单个假设的安全不能保护 ERM——因为 ERM 选中的恰恰是 \(\widehat R\) 最小、最可能``走运''的那个。但如果 \(\mathcal{H}\) 是有限集，union bound 可以把全部坏事件一次买断：

\[
P\Bigl(\bigcup_{h\in\mathcal{H}} A_h\Bigr)
\le \sum_{h\in\mathcal{H}} P(A_h)
\le 2|\mathcal{H}|\,e^{-2n\varepsilon^2}.
\]

令右端等于 \(\delta\)，解出 \(\varepsilon\)：以至少 \(1-\delta\) 的概率，

\[
|\widehat R(h) - R(h)|
\le \sqrt{\frac{\log(2|\mathcal{H}|/\delta)}{2n}}
\quad\text{对所有 } h\in\mathcal{H} \text{ 同时成立}.
\]

``同时''这个词承担了全部重量。这个不等式不是说``我挑一个假设之后再套 Hoeffding''——那是不成立的。它是说：整批候选假设的经验风险和真风险，在同一个高概率事件里全部靠近。\(\widehat R\) 作为 \(\mathcal{H}\) 上的一个函数，在这个事件里一致地逼近 \(R\)。ERM 选中谁都无所谓——它总在这个事件覆盖范围内。

代价出现在根号里。把单假设的界 \(\sqrt{\log(2/\delta)/(2n)}\) 和一致界放在一起：唯一的差别是分子从 \(\log(2/\delta)\) 变成了 \(\log(2|\mathcal{H}|/\delta)\)。多出来的 \(\log|\mathcal{H}|\) 就是选择自由度的价格。

用刚才的数字跑一遍。\(|\mathcal{H}| = 1000\)，\(n = 10000\)，\(\delta = 0.05\)：

固定假设的界宽：\(\sqrt{\log(2/0.05) / 20000} = \sqrt{3.69 / 20000} \approx 0.0136\)。

一致界的界宽：\(\sqrt{\log(2\times 1000/0.05) / 20000} = \sqrt{10.60 / 20000} \approx 0.0230\)。

为同时保护 1000 个候选，界宽从 1.36\% 扩大到 2.30\%。代价是存在的，但并不大——因为对数增长极慢。把候选数量翻倍到 2000，\(\log 2000 \approx 7.60\)，界宽增加到约 2.43\%——几乎没变。

但这是好消息的另一面。如果假设类的大小随某个``开关数'' \(k\) 指数增长——\(k\) 个二元选择给出 \(|\mathcal{H}| = 2^k\)——那么 \(\log|\mathcal{H}| = k\log 2\)，界宽按 \(\sqrt{k/n}\) 变松。每一个可以自由拨动的开关，都要花样本预算去买。

\subsection{一致收敛让 ERM 的安全有了数学凭证}

一致界一旦成立，就可以沿着一条简单的不等式链把 ERM 输出的真风险和类内最优真风险连接起来。

记 \(\widehat h \in \argmin_{h\in\mathcal{H}} \widehat R(h)\) 为 ERM 的输出，\(h^* \in \argmin_{h\in\mathcal{H}} R(h)\) 为类内真风险最优的假设。在一致界成立的好事件里：

\[
R(\widehat h)
\le \widehat R(\widehat h) + \varepsilon
\le \widehat R(h^*) + \varepsilon
\le R(h^*) + 2\varepsilon.
\]

三个不等号各司其职。第一个：一致界对 \(\widehat h\) 成立，训练误差低不意味着真风险低，但差距不超过 \(\varepsilon\)。第二个：ERM 的定义——\(\widehat h\) 在训练集上不输任何假设，包括 \(h^*\)。第三个：一致界对 \(h^*\) 也成立。

移项即得超出风险（excess risk）的界：

\[
R(\widehat h) - R(h^*)
\le 2\sqrt{\frac{\log(2|\mathcal{H}|/\delta)}{2n}}
= \sqrt{\frac{2\log(2|\mathcal{H}|/\delta)}{n}}.
\]

这个不等式给出的是相对保证：ERM 找到的假设，其真风险不比类内最优假设的真风险差超过右边这个量。它不能保证 \(R(\widehat h)\) 本身很小——如果整个 \(\mathcal{H}\) 里根本没有好假设，\(R(h^*)\) 本身就很大，一致收敛只能保证你找到了一个接近这个很差的最优值的假设。学习问题的难度同时取决于 \(\mathcal{H}\) 的表达力（能不能逼近真实规律）和样本量（能不能从有限数据中锁定好的那个）。

把不等式反过来解样本量：要以至少 \(1-\delta\) 的概率把超出风险压到 \(\varepsilon_0\) 以内，需要

\[
n \ge \frac{2\log(2|\mathcal{H}|/\delta)}{\varepsilon_0^2}.
\]

依赖结构很清晰：精度 \(1/\varepsilon_0^2\) 是二次的——想把误差再砍一半，需要四倍样本；置信度 \(\log(1/\delta)\) 是对数的——从 95\% 提到 99\% 的代价很小；类大小 \(\log|\mathcal{H}|\) 也是对数的。用 PAC（probably approximately correct）的语言说：有限假设类在 agnostic 意义下可学，且样本量仅对数依赖于类的大小。

这是统计学习理论给出的第一个正面答案。但它只控制了估计误差——由有限样本导致的那部分损失。逼近误差——\(R(h^*)\) 离问题本身的最优风险有多远——是 \(\mathcal{H}\) 的设计决定的。扩大 \(\mathcal{H}\) 会减小逼近误差，却增大 \(\log|\mathcal{H}|\)。偏差与方差的权衡已经写在这条不等式里了，本单元最后一篇会回到这个天平。

\subsection{基数的破产}

到目前为止一切都成立——对于有限 \(\mathcal{H}\)。现在把上面的公式代入任何一个现实模型，会立刻撞上同一堵墙。

\(\R\) 上的阈值分类器 \(h_t(x) = \sign(x - t)\)：阈值 \(t\) 可以取遍实数轴上的任何一个点，\(|\mathcal{H}|\) 是不可数无穷。\(\R^d\) 中的线性分类器 \(h_{w,b}(x) = \sign(w^{\trans}x + b)\)：参数空间是 \(\R^{d+1}\)，同样无穷。Union bound 对无穷多个事件求和得到``概率不超过无穷大''——这句话在数学上没有提供任何信息。

而 \hyperref[sec:09-Convex-Optimization/07-Applications/03]{``最大间隔与支持向量机''} 中那些表现良好的模型，恰恰就属于这种无限类。如果理论只能解释有限类，它就解释不了我们每天都在用的模型。

一个本能的退路是离散化。计算机里每个权重都是 64 位浮点数，\(d\) 个参数的状态总数虽然巨大但是有限，\(\log|\mathcal{H}| \approx 64d\)。代回样本量公式，确实能得到一个有限的数字。但这个答案在错的地方花了预算：同一个函数类，改用 32 位浮点表示就``便宜''了一倍，改用 128 位就``贵''了一倍。函数集合的表达能力——它能实现哪些输入-输出映射——和存储这些函数所用的比特数毫无关系。两个相差 \(10^{-9}\) 的权重向量，在任何现实精度的数据上产生完全相同的分类结果，但基数把它们计为两个不同的假设，各收一份 union bound 的税。

基数是错的计量单位。正确的单位要换一个地方去找：数据集。

无论 \(\mathcal{H}\) 中有多少个函数，把它们放到 \(n\) 个数据点上，它们能产生的标记方式最多只有 \(2^n\) 种——每个样本分两类，总共 \(n\) 个样本，最多 \(2^n\) 种排列。两个参数相差极小的线性分类器，在这 \(n\) 个点上完全可能给出相同的预测。它们在数据上的行为是同一个模式，应该被计为同一种可能。

下一节把这个观察变成定义——VC 维，`生长函数'，`打散'——并用``在有限数据上能产生多少种不同行为''替换这里的 \(\log|\mathcal{H}|\)。替换之后的界对无限类仍然有效，且不再依赖参数的存储精度。基数破产的地方，正是 VC 维开始工作的地方。
